\section{The projective-line family over the open stratum}\label{sec:projective-family}

Let $S\to T$ be the auxiliary K3 resolution, with
exceptional $(-2)$-curve $C$, and identify
$S^\circ=S\setminus C=T\setminus\{o\}$.
In the two-plane chamber write $D^\circ=D\setminus E$.
The map $\pi:D^\circ\to S^\circ$ is a projective-line
bundle. It has trivial monodromy on fibre cohomology,
yet a nonzero algebraic Brauer obstruction. This
distinction is visible in the global topology.

\subsection{Algebraic, analytic and topological lifts}
For a stable bundle $F$ of vector $w$, the groups
$\Hom(L_1,F)$ and $\Ext^2(L_1,F)$ vanish, while
$\dim\Ext^1(L_1,F)=2$. Local universal $w$-families
therefore give rank-two extension bundles. Their
projectivizations descend to $D^\circ$, with the
order-two Brauer class of the universal $w$-family.
This class remains nonzero on $S^\circ$: for a
regular integral scheme the Brauer group injects
into that of its function field
\cite[Corollaries~1.8 and~1.10]{grothendieckbrauer}.
Consequently $D^\circ$ has no presentation as the
projectivization of an ordinary \emph{algebraic}
rank-two vector bundle on $S^\circ$.

The curve $C$ is primitive in the unimodular K3
lattice, since its square is $-2$. The Gysin
sequence gives
\begin{equation}\label{eq:open-surface-cohomology}
 H^1(S^\circ)=H^3(S^\circ)=H^4(S^\circ)=0,
 \qquad H^2(S^\circ)=H^2(S)/\Z[C]\simeq\Z^{21}.
\end{equation}
The connected group $\mathrm{PGL}_2(\C)$ acts
trivially on $H^0(\PP^1)$ and $H^2(\PP^1)$.
The fibre cohomology local systems are thus
constant. There is also no topological obstruction
in degree three to lifting the projective bundle.
The class $u=\alpha|_{D^\circ}$ supplies an explicit
integral hyperplane class, because $\alpha\cdot f=1$.

An ordinary holomorphic lift exists as well. Siu's
top-degree coherent-cohomology theorem for a
connected noncompact complex manifold
\cite{siu} gives $H^2(S^\circ,\OO)=0$.
Together with $H^3(S^\circ,\Z)=0$, the exponential
sequence gives $H^2(S^\circ,\OO^*)=0$. The obstruction
from $1\to\OO^*\to\mathrm{GL}_2(\OO)
\to\mathrm{PGL}_2(\OO)\to1$ therefore vanishes.
There is no contradiction: on this nonproper
surface analytification need not retain an
algebraic Brauer obstruction.

\subsection{The divisor class, including its multiplicity}
We need the precise class of $D$ to calculate the
ring of $D^\circ$. Put $A=L_1$. The inequality
$\zeta_1<0$ means $\mu_\eta(A)>\mu_\eta(E)$,
so $\Hom(A,E)=0$ throughout the chamber.
Moreover $\chi(A,E)=-(a_1,v)=0$.

A nonzero map $E\to A$ has image
$A(-B)\otimes I_Z$ for an effective divisor $B$
and a zero-dimensional scheme $Z$. At $H$ its
image is a torsion-free quotient of a semistable
slope-zero bundle, so has nonnegative slope.
As a subsheaf of $A$ it has slope $-H\cdot B$.
Ampleness forces $B=0$. The kernel has slope zero
and is semistable. The complete factor and defect
classification of Lemma~\ref{lem:interacting-factor-list}
then forces $Z=\varnothing$ and makes the kernel
a bundle of vector $w$. Conversely the extension
fibres, including both deepest planes, have a
quotient to $A$. Hence $\Hom(E,A)\ne0$ precisely
on $D$.

The perfect complex $Rp_*R\Hom(A,\cE)$ has
amplitude $[1,2]$, virtual rank zero, and is
acyclic off $D$. Locally write it as
$K^1\xrightarrow{d}K^2$. Its determinant section
lies in $\det(K^1)^{-1}\otimes\det(K^2)$, whose
first Chern class is $\theta_v(a_1)$ by
Grothendieck--Riemann--Roch and
\eqref{eq:mukai-marking}.

We check that the section vanishes simply at a
generic point of $D$. There the bundle fits into
$0\to F\xrightarrow{j}E\xrightarrow{q}A\to0$,
with $F$ $H$-stable. The extension sequence,
$\Hom(F,A)=\Hom(A,F)=0$, and simplicity give
\[
 \begin{gathered}
 \Hom(F,E)=\C j,\quad \Hom(E,A)=\C q,\quad
 \Hom(E,E)=\C,\\
 \Ext^1(E,A)=\C,\quad
 \Ext^2(E,A)=0.
 \end{gathered}
\]
In the exact sequence
\[
 \Ext^1(E,E)\xrightarrow{q_*}\Ext^1(E,A)
 \longrightarrow\Ext^2(E,F)
 \xrightarrow{j_*}\Ext^2(E,E),
\]
the last arrow is an isomorphism: its Serre dual
sends $\operatorname{id}_E$ to $j$ between
one-dimensional spaces. Thus $q_*$ is onto.
For an infinitesimal deformation $\delta$ of $E$,
$q_*\delta$ is the obstruction to extending $q$.
Serre duality identifies this obstruction with
the first variation of the local differential
from $\Ext^1(A,E)$ to $\Ext^2(A,E)$, up to sign.
After cancelling acyclic summands, $d$ is a scalar
function with nonzero derivative. Its zero is
simple. These tangents preserve the determinant
because $H^1(X,\OO_X)=0$.

Since $D$ is irreducible and is the whole
divisorial support, we have proved
\begin{equation}\label{eq:divisor-class}
 [D]=\theta_v(a_1),\qquad D|_{P_i}=-h_i,
 \qquad D\cdot f=-2.
\end{equation}
The multiplicity argument is needed: the last
two intersection numbers by themselves would
not identify the global divisor class.

\subsection{The ring over the open stratum}
Use the integral isotropic Mukai map
$\theta_w:w^\perp/\Z w\to H^2(S,\Z)$ induced
by the same $w$-family as in the construction.
With $x$ from \eqref{eq:restriction-witnesses}, set
\[
 z=a_1+2x,\qquad d=\theta_w(z)|_{S^\circ}.
\]
Direct pairing gives $z\in v^\perp$,
$(z,a_1)=0$ and $(z,w)=0$. The universal
extension over $D^\circ$ gives
\begin{equation}\label{eq:mukai-open-compatibility}
 \theta_v(z)|_{D^\circ}=\pi^*\theta_w(z)|_{S^\circ}.
\end{equation}
Here is the descent argument for this global
identity. On the universal gerbe, use the
normalized rational Chern character of the
twisted family; equivalently divide the Chern
character of a quasi-universal family by its
similitude. Additivity in the universal extension
gives \eqref{eq:mukai-open-compatibility} in rational
cohomology. A tautological line, a base line or
a rational choice for the torsion gerbe affects
the degree-two pushforward only through
$(z,a_1)$, $(z,w)$ or $(z,v)$, all zero.
Thus the identity descends globally, rather
than only on contractible charts. Both sides
are integral. Leray--Hirsch with $u$ and
\eqref{eq:open-surface-cohomology} shows that
$H^2(D^\circ,\Z)$ is torsion-free, so the
rational equality is integral.

Combining it with \eqref{eq:divisor-class} gives
\begin{equation}\label{eq:relative-canonical}
 D|_{D^\circ}=-2u+\pi^*d.
\end{equation}
By adjunction this is the relative canonical
class: $M$ and $S^\circ$ have trivial canonical
bundles. Choose a topological rank-two lift $\mathcal W$
of the projective bundle, in the quotient
convention for $\PP(\mathcal W)$. Its hyperplane class
differs from $u$ by a pullback from the base.
Tensoring $\mathcal W$ with a topological line bundle
removes this difference. The normalized lift
has relative canonical class
$-2u+\pi^*c_1(\mathcal W)$, so injectivity of $\pi^*$
gives $c_1(\mathcal W)=d$. Also $c_2(\mathcal W)=0$ by
\eqref{eq:open-surface-cohomology}. The integral
projective-bundle formula now yields
\begin{equation}\label{eq:open-divisor-ring}
 H^*(D^\circ,\Z)
   =H^*(S^\circ,\Z)[u]/(u^2-du).
\end{equation}
This calculation specifies the hyperplane
normalization. It is consistent with the
nonzero algebraic Brauer class and the trivial
fibre-cohomology monodromy: these record
different aspects of the same family.
